\begin{papercard}{kelley2009}
\cf{Stages} \stagetag{C} (also \stagetag{S} \stagetag{W}). An iterative design study of a standardised ``Privacy Nutrition Label''. It moves through four prototypes and two focus groups, then a within-subjects lab study ($n=24$) of accuracy, time and enjoyment against natural-language policies. It is the conceptual ancestor of today's app-store privacy labels (Apple, Google Data Safety).

\cf{Problem} Policies are hard to understand (jargon, college reading level, vagueness), people do not read them, and consumers believe they have no choices \pc{kelley2009}{1}{consumers do not actually believe they have choices when it comes to their privacy}. Goal: a clear, uniform, single-page summary that fixes three problems: finding information, recognising differences between policies, and reading time \pc{kelley2009}{1}{We set out to design a clear, uniform, single-page summary}.

\cf{Design method} A review of other labelling regimes: Nutrition Facts, energy, drug facts, the FTC/Kleimann financial notice and traffic-light labels \pc{kelley2009}{3}{We also explored energy labeling programs from the European Union}. Iteration: P3P Expandable Grid, then Simplified Label, then Simplified Grid, then Privacy Nutrition Label. Five problems with the earlier grid were identified, e.g.\ 40.7\,\% of users never expanded any rows \pc{kelley2009}{4}{many users (40.7\%) never expanded any data types}. Principles from nutrition labels: a bounding box, bold rules and a clear title \pc{kelley2009}{4}{Putting a box around the label identifies the boundaries of the information}. Data model: P3P triplets \{data, purpose, recipient\}. Retention was dropped because most policies say ``indefinitely'' \pc{kelley2009}{4}{A useful model is to think of P3P as consisting of multiple triplets}. Flattening to a list loses information and makes policies look more permissive \pc{kelley2009}{5}{This tends to make privacy policies appear more permissive than they actually are}.

\cf{Design lessons from testing} Blank cells were read as ``anything goes'', so an explicit ``not collected'' symbol was needed \pc{kelley2009}{7}{in testing, many participants were unclear on the meaning of the blank cells}. All ten data categories are always shown in a fixed order, so labels can be compared side by side \pc{kelley2009}{7}{Any policy displayed in this format will have exactly 10 rows}. The ``mixed choices'' symbol was incomprehensible and was replaced by showing the most invasive practice \pc{kelley2009}{7}{Instead we now display the symbol for the most invasive practice}. Darker means worse, colours are colour-blind-safe, and a glossary of 17 ``useful terms'' was added \pc{kelley2009}{7}{consists of seventeen definitions, one for each of the row and column headers}.

\cf{Focus groups} Two one-hour groups ($n=10$, $n=7$). Participants could spot differences but in the first group half chose the more invasive company. Opt-in and opt-out were poorly understood \pc{kelley2009}{8}{five of the ten participants chose Policy B}.

\cf{Lab study} Within subjects, with random order of label vs.\ natural language (NL), and the same 24 questions per format \pc{kelley2009}{8}{We used a within-subjects design where participants were randomly assigned}. Stimuli: a real e-commerce P3P policy modified on three dimensions (collection, profiling, opt-in vs.\ opt-out sharing) to create fictional companies with a clear correct answer. The NL policies were matched in length and readability (FRES about 42, FKGL about 11.6) \pc{kelley2009}{8}{These significant differences were introduced so that there would be a clear}, \pc{kelley2009}{9}{modified the second policy so that both were approximately the same length}. Tasks: 8 information-finding questions, 6 Likert items, 6 comparison items, 4 enjoyment items \pc{kelley2009}{9}{The task structure for each condition was exactly the same, with 24 tasks}.

\cf{Analysis and results} Correct answers: label mean 10.13 vs.\ NL 6.83 ($t(23)=7.41$, $p<.001$) \pc{kelley2009}{10}{We compared the correct number of total questions, per participant}. Per-question McNemar tests with a Benjamini--Hochberg correction for multiple testing \pc{kelley2009}{10}{We performed a Benjamini-Hochberg correction to account for multiple testing}. Time: 339.9\,s vs.\ 692.0\,s in total, with timing analysed only for correct answers. Satisfaction was significantly better on 8 of 10 items. Order effects (priming) were present: the label was rated higher by those who saw NL first \pc{kelley2009}{10}{These results show significant priming to appreciate the grid more}.

\cf{Limitations acknowledged} Paper-based lab tasks, a small student-heavy sample, and open questions about real-world use over time \pc{kelley2009}{11}{A number of open questions remain about how people will use the label in practice}.

\cf{Use for this thesis} Three uses. (1) The \emph{data--purpose--recipient triplet} is a compact coding unit for policies and labels, and the warning that flattening it overstates or blurs practices applies to how Data Safety sections are coded. (2) Design principles (fixed rows, explicit ``not collected'', no ambiguous mixed symbols) give criteria for judging app-store labels. (3) Methodologically: a within-subjects comparison with counterbalancing, McNemar tests with multiple-testing correction, and report order effects. Treat the 2009 P3P context as historical. It is a lab study of a prototype, not evidence about current store labels (see Li et al.\ 2022).
\end{papercard}
